\documentclass{jsarticle}
\usepackage[dvipdfmx, final]{graphicx}
\usepackage{amsmath}
\usepackage{amssymb}
\begin{document}
\title{M theory equal with AdS5 manifold.}
\author{}
\date{}
\maketitle

   \newcommand{\variint}{%
	\begin{picture}(15,30)
		\put(0,0){
			$ \iint $}
		\put(0,5){\line(15,0){15}}
	\end{picture} %
}

\newcommand{\variiint}{%
	\begin{picture}(15,15)
		\put(0,0){
			$ \iiint $}
		\put(0,5){\line(15,0){20}}
	\end{picture} %
}



 \newcommand{\alearletter}{%
	 \begin{picture}(10,20)
		 \put(0,0){
			 $ \square $}
		 \put(0,0){\line(1,1){10}}
	 \end{picture} %
	 }


       ブライアン・グリーン教授の宇宙からの演繹された説明で、
    場の理論となるオイラーのβ関数を背景にして、
    ゼータ関数が単体となり、それを元に構築された素粒子が
    Jones多項式として、各方程式の原本になり、
    それと同型の特殊相対性理論の多様体積分が、
    一般相対性理論の多様体積分として、ガンマ関数の大域的積分多様体
    として、オイラーの判別式のノルムのD-braneの種数として、
    M理論としてのカラビ・ヤウ多様体が、AdS5多様体と同型と言えるのを、
    始めは、特殊相対性理論と一般相対性理論の説明をされて、
    それから、量子力学の不確定性原理を述べて、
    そのデータが、オイラーのβ関数を場の理論とすると、
    全部の重力場と量子力学レベルの素粒子方程式が、
    同型となる、Jones多項式へと連想される説明で、
    相棒のリサ・ランドール博士と本当に友人なのですね、
    との感想で、このthe elegant Universeの本を読んでの
    感想です。
    表紙が、種数が3の多様体を描いていて、サーストン・ペレルマン多様体
    の1部分空間を示していて、この種数3が$E^{0}\times S^{2} $と
    長生きする宇宙と、エロさまで表している、ブライアン・グリーン教授の大人の
    表現に、いいなと思いました。

    11次元多様体は、10次元で重力場を表して、11次元目でディラトンが
    出てきて、それが、AdS5多様体では、4次元目で重力場と反重力場を
    表していて、5次元目に電磁場を表していて、
    この11次元と5次元が、種数で同型と言えて、両方が、M理論を
    別の側面から見ているとリサ・ランドール博士とブライアン・グリーン教授
    は、過去も未来も現在も述べている。

    種数1の代数幾何の量子化にm,nを加群した代数幾何の量子化の加群同士で積としての、
    環を求めると、ベータ関数での種数3の多様体となる。これは、サーストン・ペレルマン多様体の一部である
    幾何構造であり、the elegant universeの本の表紙を表している。綺麗な宇宙である。
         代数的計算手法のために$ \oplus L $を使っている。
そのために、冪乗計算と商代数の計算が、乗算で楽に見えるようになっている。
            微分幾何の量子化は、代数幾何の量子化の計算になっている。
  加減乗除が初等幾何であり、大域的微分と積分が、現代数学の代数計算の
  簡易での楽になる計算になっている。
  初等代数の計算は、
  $$ \bigoplus({i \hbar^{\nabla}})^{\oplus L} $$
   $$m \bigoplus({i \hbar^{\nabla}})^{\oplus L} +  
   n\bigoplus({i \hbar^{\nabla}})^{\oplus L} = 
   (m+n) \bigoplus({i \hbar^{\nabla}})^{\oplus L}   $$
   $$ m\bigoplus({i \hbar^{\nabla}})^{\oplus L} -  
   n\bigoplus({i \hbar^{\nabla}})^{\oplus L} = 
   (m-n) \bigoplus({i \hbar^{\nabla}})^{\oplus L}   $$
   $$ \bigoplus({i \hbar^{\nabla}})^{{\oplus L}^m} \times 
      \bigoplus({i \hbar^{\nabla}})^{{\oplus L}^n} =
   \bigoplus({i \hbar^{\nabla}})^{\oplus L^{m+n}} $$
  $$ {{\bigoplus({i \hbar^{\nabla}})^{\oplus L^m}}  \over  
  {\bigoplus({i \hbar^{\nabla}})^{\oplus L^n}}} =  
  \bigoplus({i \hbar^{\nabla}})^{\oplus L^{m\over n}} $$
  大域的計算での微分と積分は、
  $$ \left(\bigoplus({i \hbar^{\nabla}})^{\oplus L}\right)^{df}
  =  {\left(\bigoplus({i \hbar^{\nabla}})^{\oplus L}\right)^{
	  \bigoplus({i \hbar^{\nabla}})^{\oplus L}}}^{'}  $$
  $$ = \bigoplus({i \hbar^{\nabla}})^{\oplus L^{'}} $$
   $$ \int \bigoplus({i \hbar^{\nabla}})^{\oplus L} dx_m $$
   $$ \bigoplus({i \hbar^{\nabla}})^{\oplus L} = n $$
   大域的積分の代数幾何の量子化での計算は、ガンマ関数になっている。
   $$ {n^{L+1} \over {n+1}} = \int e^x x^{1 - t} dx $$
   代数幾何の量子化の因数分解は、加減乗除の式のm,nの
   組み合わせ多様体でのガンマ関数同士の計算になる。


       This estern with Gamma function resteamed from being riging to Beta function in Thurston
   Perelman manifold. This field call all of theorem to architect with Space ideal of quantum level.
     This theorem will be estern the man to be birth with Japanese person.
   This person pray with be birth of my son.
   This pray call work to be being name to say me pray.
   This pray resteam me to masterbation and this play realized me Gakkari.
   Aya san kill me to be played.

     I like this poem to proof with English moreover Japanese language loved from me.
     And this crystal proof released me to write English and Japanese language to discover them
     from mathmatics theorems.

   $$ \left(\bigoplus({i \hbar^{\nabla}})^{\oplus L}+m\right) 
   \left(\bigoplus({i \hbar^{\nabla}})^{\oplus L}+n\right) $$
  $$ = {n^{L+1} \over {n+1}} = t \int (1 - x)^n x^{m - 1} dx $$



     This equation esterminate with Beta function in Gamma function riginged from telphone
  to world line surface. And this ringed have with Algebra manifold of differential geometry in
  quantum level. This write in English language. Moreover that' cat call them to birth of Japanese cats.
  And moreover, I birth to name with Japanese Person. 
  And, this theorem certicefate the man to birth Diths Person.
  This stimeat with our constrate with non relate person and cat.


  $$ = \int x^{m-1}(1 - x)^{n-1} dx $$


  $$ \nabla ({i \hbar^{\nabla}})^{\oplus L}, \bigoplus ({i \hbar^{\nabla}})^{
	   \oplus L}, \square ({i \hbar^{\nabla}})^{\oplus L} $$
   $$ \boxtimes (i \hbar^{\nabla})|_{dx_m}^L, \boxplus (i \hbar^{\nabla}|_
   {dx_m}^L) $$



  $$ x^{{1 \over 2} + iy} = e^{x \log x} $$
  それゆえに、この定数はゼータ関数と微分幾何の量子化を因数にもつ
  素因子分解の式にもなっている。
  Therefore, this product also constructed with differential geometry
  of quantum level equation and zeta function.
  $$ = C $$ 
  そして、この関数はオイラーの定数から広中平祐定理による4重帰納法の
  オイラーの公式からの多様体積分へとのサーストン空間のスペクトラム関数
  ともなっている。
  

       大域的多重積分と大域的無限微分多様体を定義すると、
   Global assemble manifold defined with infinity of differential
   fields to determine of definition create from Euler product and
   Gamma function for Beta function being potten from integral manifold
   system. Therefore, this defined from global assemble manifold from
   being Gamma of global equation in topological expalanations.

   $$ \int f(x)dx = \int \Gamma(\gamma)^{'}dx_m $$
   $$ = 2 (\cos (i x \log x) - i \sin(i x \log x)) $$
   $$ \left({{\int f(x) dx} \over \log x}\right) = \lim_{\theta \to \infty}
    \left({{\int f(x) dx} \over \theta}\right) = 0,1 $$
    $$ e^{i \theta} = \cos \theta + i \sin \theta $$
    $$ \left({\int f(x) dx}\right)^{'} = 2 (i \sin(i x \log x) - \cos (i x \log
    x)) $$
    $$ = 2(- \cos(i x \log x) + i \sin(i x \log x)) $$
    $$ (\cos (i x \log x) - i \sin(i x \log x))^{'} $$
    $$ = {d \over {d {e^{i\theta}}}}\left((\cos, - \sin )\cdot
    (\sin, \cos)\right)  $$
    $$ \int \Gamma(\gamma)^{''}dx_m = \left({\int \Gamma(\gamma)^{'}}dx_m
    \right)^{\nabla L} = \left({{\int \Gamma dx_m} \cdot {{d \over d\gamma}
    \Gamma}}\right)^{\nabla L} \le 
    \left({\int \Gamma dx_m + {d \over d\gamma}\Gamma}\right)^{\nabla L} \le 
    \left({e^{f} - e^{-f}} \le {e^{-f} + e^{f}}\right)^{'}  $$
    $$ = 0,1 $$

    Then, the defined of global integral and differential manifold
    from being assembled world lines and partial equation of deprivate
    formula in Homology manifold and gamma function of global topology
    system, this system defined with Shwaltsshild cicle of norm from
    black hole of entropy exsented with result of monotonicity for
    constant of equations.
    以上 より、大域的微分多様体を大域的 2重微分多様体として、
    処理すると、ホモロジー多様体では、種数が１であり、特異点では、種数が０
    と計算されることになる。ガンマ関数の大域的微分多様体では、
    シュバルツシルト半径として計算されうるが、これを大域的2重微分で
    処理すると、ブラックホールの特異点としての解が無になる。




	Abel拡大 $ K/k $に対して、
	$$ f = {\pi}_{p}{f_p} $$
        類体論
	Artin記号を用いて、
       $$ \left({{\alpha, K/k} \over {p}}\right) 
       = \left({{K/k} \over {b}}\right) (\in G) $$
       $ {\alpha/ \alpha_0} \equiv 1 \pmod {f_p} $, 
      $ \alpha_0 \equiv 1 \pmod {f f_p^{-1}} \to  
      \alpha \in k $
      $ (\alpha_0) = {p}^{\alpha}b, \text{
	      $p$ と$b$は互いに素} $
      $b \to \text{相対判別式$\delta K/k$で互いに素}$
      この値は、補助数${\alpha_0}$の値の取り方によらずに、
      一意的に定まる。
      $$ \left({{\alpha, K/k} \over {p_{\infty}^{(j)}}}\right)  = 
      \text{1または0} $$

      これらをまとめた式が、Hilbertの剰余記号の判別式
      $$ \mathcal{\pi}_{p}\left({{\alpha,b} \over 
      {p}}\right) = 1 $$
      であり、
    　この式たちから、代数幾何の種数のノルム記号である、
$$ ||ds^2|| = \lim_{x \to \infty} [\delta(x) \int \int \int \pi \left( \sum_{k=0}^{\infty}
		{{}^n\sqrt{p},x \over n} \right)
^{1 \over 2}d \tau]^{\mu\nu} $$
     が求まり、
     $$ p^{\alpha} n = {}^n \sqrt{p} $$
     $$ n^{{}^n\sqrt{p}} = \bigoplus{(i \hbar^{{{\nabla)}}^{\oplus 
     L}}} $$
     $$ = n^{p^{{1 \over n}}} = n^{{-n}^{p}}  $$
     $$ = \int \Gamma(\gamma)^{'}dx_m = e^{-x \log x} $$
     となり、
     kの素イデアルの密度Mに対して、
     $$ \lim_{s \to {1 + 0}} \sum_{{p \in M}} {
	     1 \over (N(p))^{s}}/ \log {1 \over 
     {s - 1} } $$
     $$ = \text{Mの密度(density)} $$

     $$ \alpha(f{d \over dt},g{d \over dt}) = \int_s \begin{vmatrix}
	     f^{'} & f^{''} \\
	     g^{'} & g^{''} \end{vmatrix}dt,
	     \mathcal{B}(f{d \over dt},g{d \over dt},h{d \over dt}) =
	     \int_s \begin{vmatrix} f & f^{'} & f^{''} \\
		     g & g^{'} & g^{''} \\
	     h & h^{'} & h^{''} \end{vmatrix}dt $$
    これらは、Gul'faid-Fuksコホモロジーの概念を局所化することにより、
    形式的ベクトル和のつくるコホモロジーとして、導かれている。



        代数幾何の量子化では、種数1であり、閉３次元多様体では、種数２であり、
    ガンマ関数の和と積の商代数では、ベータ関数として、種数０であり、
    ランクから、代数幾何の量子化の加群同士では、代数幾何の量子化が、ワームホールを
    種数１持っていて、この加群で、係数tのベータ関数となり、種数３のワームホール２種の
    ベータ関数となっている。これを整理すると、閉３次元多様体にワームホール１種が加わっている
    ベータ関数が$E^{0} \times S^{2} $と、種数１のベータ関数に２種のワームホールがあり、
    合計種数が３種の代数幾何になっている。


      これが、the elegant universeの表紙に載っている図になっている。

      種数０の補空間が種数１であり、種数１の補空間が種数２であり、
      種数２の補空間が種数３である。

      時間の一方向性が、電磁場理論の電弱相互理論であり、時間が電磁場である。
      １１次元多様体の１０次元が重力で、１１次元目が電磁場、ディラトンが時間である。
      これは、種数が３であり、５次元多様体の種数が３と同型である。
      ３次元多様体が種数が２である。
      これにワームホール１種であり、種数が３になる。
      表裏が表裏一体になっている。

      代数幾何の量子化の加群同士でも、ベータ関数となり、種数が３になる。
      ウィッテンが１１次元超重力理論を提出していることを、
      $$ e^{-x \log x} \le y \le e^{x \log x}, y \neq 0 $$
      と、フェルマーの定理の解を範囲に値をとる。

      すべては、Jones多項式が統一場理論となる。

$$ ||ds^2|| = \lim_{x \to \infty} [\delta(x) \int \int \int \pi \left( \sum_{k=0}^{\infty}
		{{}^n\sqrt{p},x \over n} \right)
^{1 \over 2}d \tau]^{\mu\nu} $$
$$ \pi \left( \sum_{k=0}^{\infty}
		{{}^n\sqrt{p},x \over n} \right)
^{1 \over 2}d \tau]^{\mu\nu} = e^{-f}dV $$

     が求まり、
     $$ p^{\alpha} n = {}^n \sqrt{p} $$
     $$ n^{{}^n\sqrt{p}} = \bigoplus{(i \hbar^{{{\nabla)}}^{\oplus 
     L}}} $$
     $$ = n^{p^{{1 \over n}}} = n^{{-n}^{p}}  $$
     $$ = \int \Gamma(\gamma)^{'}dx_m = e^{-x \log x} $$
     となり、
     kの素イデアルの密度Mに対して、
     $$ \lim_{s \to {1 + 0}} \sum_{{p \in M}} {
	     1 \over (N(p))^{s}}/ \log {1 \over 
     {s - 1} } $$
     $$ = \text{Mの密度(density)} $$

     $$ \alpha(f{d \over dt},g{d \over dt}) = \int_s \begin{vmatrix}
	     f^{'} & f^{''} \\
	     g^{'} & g^{''} \end{vmatrix}dt,
	     \mathcal{B}(f{d \over dt},g{d \over dt},h{d \over dt}) =
	     \int_s \begin{vmatrix} f & f^{'} & f^{''} \\
		     g & g^{'} & g^{''} \\
	     h & h^{'} & h^{''} \end{vmatrix}dt $$
    これらは、Gul'faid-Fuksコホモロジーの概念を局所化することにより、
    形式的ベクトル和のつくるコホモロジーとして、導かれている。


    $$ p = e^{x \log x},e^{-x \log x} $$
    $$ p = \bigoplus{(i \hbar^{\nabla})}^{\oplus L} $$
    pの取り得る範囲で、Hilbert多様体は、
    $$ ||ds^2|| = 0,1 $$
    の種数の値を取る。
    この補空間が種数３である。

\end{document}
